\documentclass[10pt]{beamer}
\usetheme{metropolis}
\usepackage{graphicx}
\usepackage{listings}
\usepackage{booktabs}
\usepackage{xcolor}
\usepackage{hyperref}

\definecolor{codebg}{HTML}{F5F5F5}
\definecolor{codegreen}{HTML}{2E7D32}
\definecolor{codegray}{HTML}{757575}
\definecolor{codeblue}{HTML}{1565C0}

\lstdefinestyle{rstyle}{
  language=R, backgroundcolor=\color{codebg},
  basicstyle=\ttfamily\scriptsize,
  keywordstyle=\color{codeblue}\bfseries,
  stringstyle=\color{codegreen},
  commentstyle=\color{codegray}\itshape,
  breaklines=true, frame=single,
  rulecolor=\color{codegray!30}, showstringspaces=false, tabsize=2,
  morekeywords={library,ggplot,aes,geom_line,geom_text,geom_col,
    theme,element_text,element_blank,labs,annotate,ggsave,
    plot_annotation,patchwork,coord_flip,margin}
}
\lstdefinestyle{pythonstyle}{
  language=Python, backgroundcolor=\color{codebg},
  basicstyle=\ttfamily\scriptsize,
  keywordstyle=\color{codeblue}\bfseries,
  stringstyle=\color{codegreen},
  commentstyle=\color{codegray}\itshape,
  breaklines=true, frame=single,
  rulecolor=\color{codegray!30}, showstringspaces=false, tabsize=4,
  morekeywords={import,from,as,pd,plt,sns,np,fig,ax,GridSpec}
}

\title{Typography, Layout \& Composition}
\subtitle{Workshop 1 --- Module 5}
\author{Prof.Asoc.\ Endri Raco}
\institute{Data Visualization for Data Scientists \\ UNYT Tirana}
\date{2025/26}

\begin{document}

\begin{frame}
  \titlepage
\end{frame}

% ================================================================
\begin{frame}{Module Roadmap}
  \begin{enumerate}
    \item Typographic hierarchy: title, subtitle, body, caption
    \item Font selection for data visualization
    \item Whitespace, margins, and breathing room
    \item Grid systems and multi-panel composition
    \item Direct labelling, annotation, and export formats
  \end{enumerate}
  \begin{exampleblock}{Learning Outcome}
    You will be able to construct professional chart layouts with
    clear typographic hierarchy, purposeful whitespace, and
    publication-ready export settings.
  \end{exampleblock}
\end{frame}

% ================================================================
\section{Typographic Hierarchy}
% ================================================================

\begin{frame}{The Five Levels of Chart Typography}
  \begin{center}
    \includegraphics[width=0.85\textwidth]{figures_m05/type_hierarchy}
  \end{center}
\end{frame}

% ================================================================
\begin{frame}{Sizing Rules}
  \centering
  \begin{tabular}{llrl}
    \toprule
    \textbf{Element} & \textbf{Weight} & \textbf{Size} & \textbf{Purpose} \\
    \midrule
    Title    & Bold    & 16--20 pt & The ``headline'' of the chart \\
    Subtitle & Regular & 11--13 pt & Context, time range, units \\
    Axis labels & Regular & 10--11 pt & Variable names \\
    Tick labels & Regular & 8--10 pt  & Scale values \\
    Annotation  & Regular & 8--9 pt  & Direct labels on data \\
    Caption  & Italic  & 7--8 pt  & Source, notes, caveats \\
    \bottomrule
  \end{tabular}

  \vspace{6pt}
  \begin{alertblock}{Pro-Tip}
    Maintain a \textbf{minimum 1.2$\times$ ratio} between adjacent
    hierarchy levels. If your title is 16pt, the subtitle should be
    $\leq$ 13pt and axis labels $\leq$ 11pt.
  \end{alertblock}
\end{frame}

% ================================================================
\begin{frame}{Font Selection}
  \begin{center}
    \includegraphics[width=0.92\textwidth]{figures_m05/font_pairing}
  \end{center}
  \begin{exampleblock}{Data Insight}
    Sans-serif typefaces (Fira Sans, Helvetica, Arial) render crisply at
    small sizes on screens and are the default choice for all chart text.
    Reserve serif for reports; monospace for code and aligned numbers.
  \end{exampleblock}
\end{frame}

% ================================================================
\begin{frame}[fragile]{Setting Fonts in R}
  \begin{columns}[T]
    \begin{column}{0.52\textwidth}
\begin{lstlisting}[style=rstyle]
library(ggplot2)

# Global font via theme
ggplot(df, aes(x, y)) +
  geom_point() +
  theme_minimal(
    base_size = 11,
    base_family = "Fira Sans"
  ) +
  theme(
    plot.title = element_text(
      size = 16,
      face = "bold",
      margin = margin(b = 8)),
    plot.subtitle = element_text(
      size = 11,
      color = "grey50"),
    plot.caption = element_text(
      size = 8,
      face = "italic",
      hjust = 0)
  ) +
  labs(
    title = "Main Title",
    subtitle = "Context line",
    caption = "Source: ...")
\end{lstlisting}
    \end{column}
    \begin{column}{0.45\textwidth}
      \textbf{R font workflow:}

      \begin{itemize}
        \item \texttt{base\_family} sets the
              default for all text
        \item \texttt{element\_text()} overrides
              individual elements
        \item \texttt{showtext::font\_add\_google()}
              loads Google Fonts
        \item \texttt{ragg} package renders
              custom fonts in PNG/PDF
      \end{itemize}
    \end{column}
  \end{columns}
\end{frame}

% ================================================================
\begin{frame}[fragile]{Setting Fonts in Python}
  \begin{columns}[T]
    \begin{column}{0.52\textwidth}
\begin{lstlisting}[style=pythonstyle]
import matplotlib.pyplot as plt

# Global rcParams
plt.rcParams.update({
    "font.family": "sans-serif",
    "font.sans-serif": [
        "Fira Sans", "Helvetica",
        "Arial"],
    "font.size": 11,
    "axes.titlesize": 14,
    "axes.titleweight": "bold",
    "axes.labelsize": 11,
    "xtick.labelsize": 9,
    "ytick.labelsize": 9,
})

fig, ax = plt.subplots()
ax.set_title("Main Title",
    fontsize=16, fontweight="bold")
# Subtitle via fig.text or
# ax.text below title
fig.text(0.5, 0.92,
    "Context line",
    ha="center", fontsize=11,
    color="grey")
\end{lstlisting}
    \end{column}
    \begin{column}{0.45\textwidth}
      \textbf{Python font workflow:}

      \begin{itemize}
        \item \texttt{rcParams} sets global
              defaults for all plots
        \item \texttt{fontweight}, \texttt{fontsize}
              override per-element
        \item For Google Fonts: download
              \texttt{.ttf} and register via
              \texttt{font\_manager.fontManager.addfont()}
        \item \texttt{fig.text()} places text in
              figure coordinates
      \end{itemize}
    \end{column}
  \end{columns}
\end{frame}

% ================================================================
\section{Whitespace \& Margins}
% ================================================================

\begin{frame}{Whitespace Is Structure, Not Waste}
  \begin{center}
    \includegraphics[width=0.92\textwidth]{figures_m05/whitespace}
  \end{center}
\end{frame}

% ================================================================
\begin{frame}{Margin Guidelines}
  \begin{columns}[T]
    \begin{column}{0.48\textwidth}
      \textbf{Internal margins:}
      \begin{itemize}
        \item Title $\rightarrow$ plot area: 8--12 pt
        \item Plot area $\rightarrow$ axis labels: 6--8 pt
        \item Between facet panels: 8--10 pt
        \item Legend $\rightarrow$ plot edge: 6 pt
      \end{itemize}
    \end{column}
    \begin{column}{0.48\textwidth}
      \textbf{External margins:}
      \begin{itemize}
        \item Document margin: 0.5--1.0 inch
        \item Chart inset from text: 12--18 pt
        \item Between stacked charts: 18--24 pt
        \item Caption below chart: 4--6 pt
      \end{itemize}
    \end{column}
  \end{columns}

  \vspace{6pt}
  \begin{alertblock}{Pro-Tip}
    In R use \texttt{theme(plot.margin = margin(t, r, b, l))}. In Python
    use \texttt{plt.tight\_layout(pad=2.0)} or \texttt{plt.subplots\_adjust()}.
  \end{alertblock}
\end{frame}

% ================================================================
\section{Alignment \& Direct Labelling}
% ================================================================

\begin{frame}{Alignment: Numbers Right-Align}
  \begin{center}
    \includegraphics[width=0.92\textwidth]{figures_m05/alignment}
  \end{center}
  \begin{exampleblock}{Data Insight}
    Right-aligning numeric labels on the same column ensures that the
    digits align vertically, making magnitude comparison effortless.
    Centred numbers scatter the digits, forcing serial reading.
  \end{exampleblock}
\end{frame}

% ================================================================
\begin{frame}{Direct Labelling vs Legend}
  \begin{center}
    \includegraphics[width=0.92\textwidth]{figures_m05/direct_label}
  \end{center}
  \begin{alertblock}{Pro-Tip}
    A legend forces the eye to shuttle between the chart and the legend
    box. Direct labels at line endpoints (or on bars) eliminate this
    cognitive overhead. Use \texttt{ggrepel} in R or \texttt{adjustText}
    in Python to prevent label overlap.
  \end{alertblock}
\end{frame}

% ================================================================
\begin{frame}[fragile]{Direct Labelling in R}
  \begin{columns}[T]
    \begin{column}{0.52\textwidth}
\begin{lstlisting}[style=rstyle]
library(ggplot2)
library(ggrepel)

ggplot(df, aes(x = month,
               y = revenue,
               color = product)) +
  geom_line(linewidth = 1) +

  # Direct label at last point
  geom_text_repel(
    data = df |>
      filter(month == max(month)),
    aes(label = product),
    nudge_x = 0.5,
    size = 3,
    fontface = "bold",
    direction = "y") +

  # Remove legend (redundant)
  theme_minimal() +
  theme(legend.position = "none") +
  labs(x = "Month",
       y = "Revenue ($K)")
\end{lstlisting}
    \end{column}
    \begin{column}{0.45\textwidth}
      \textbf{\texttt{ggrepel} features:}

      \begin{itemize}
        \item Automatic collision avoidance
        \item Connector lines from label to
              data point
        \item \texttt{direction = "y"} constrains
              repulsion to vertical axis
        \item Works with \texttt{geom\_text\_repel}
              and \texttt{geom\_label\_repel}
      \end{itemize}
    \end{column}
  \end{columns}
\end{frame}

% ================================================================
\begin{frame}[fragile]{Direct Labelling in Python}
  \begin{columns}[T]
    \begin{column}{0.52\textwidth}
\begin{lstlisting}[style=pythonstyle]
import matplotlib.pyplot as plt
# pip install adjustText
from adjustText import adjust_text

fig, ax = plt.subplots()

for name, data in series.items():
    ax.plot(months, data,
        color=colors[name], lw=2)

    # Direct label at endpoint
    ax.text(
        months[-1] + 0.3,
        data[-1],
        name,
        fontsize=8,
        fontweight="bold",
        color=colors[name],
        va="center")

# Auto-adjust overlapping labels
# texts = [ax.text(...) for ...]
# adjust_text(texts)

ax.set_xlim(1, max(months) + 3)
\end{lstlisting}
    \end{column}
    \begin{column}{0.45\textwidth}
      \textbf{Python pattern:}

      \begin{enumerate}
        \item Place labels at line endpoints
              via \texttt{ax.text()}
        \item Extend \texttt{xlim} to accommodate
              label width
        \item Use \texttt{adjustText} for
              automatic collision avoidance
        \item Match label colour to line colour
              for visual linking
      \end{enumerate}
    \end{column}
  \end{columns}
\end{frame}

% ================================================================
\section{Annotation Patterns}
% ================================================================

\begin{frame}{Three Annotation Patterns}
  \begin{center}
    \includegraphics[width=0.92\textwidth]{figures_m05/annotation_patterns}
  \end{center}
  \begin{itemize}
    \item \textbf{Callout}: arrow + text box pointing to a specific event
    \item \textbf{Context band}: shaded region highlighting a time period
    \item \textbf{Reference line}: dashed horizontal for mean / target / benchmark
  \end{itemize}
\end{frame}

% ================================================================
\section{Grid Systems \& Composition}
% ================================================================

\begin{frame}{Layout Patterns}
  \begin{center}
    \includegraphics[width=0.92\textwidth]{figures_m05/grid_layouts}
  \end{center}
  \begin{exampleblock}{Data Insight}
    Single-chart layouts suit presentation slides. Two-column layouts
    support comparison (before/after, R vs Python). Dashboard grids
    (KPI row + chart grid) are standard in analytical reports.
  \end{exampleblock}
\end{frame}

% ================================================================
\begin{frame}{Multi-Panel Composition}
  \begin{center}
    \includegraphics[width=0.85\textwidth]{figures_m05/multipanel}
  \end{center}
\end{frame}

% ================================================================
\begin{frame}[fragile]{Multi-Panel in R: patchwork}
  \begin{columns}[T]
    \begin{column}{0.52\textwidth}
\begin{lstlisting}[style=rstyle]
library(patchwork)

# Compose with operators
# | = side by side
# / = stacked
layout <- (p1 | p2 | p3) / p4

# Annotate the composite
layout +
  plot_annotation(
    title = "Quarterly Report",
    subtitle = "Revenue metrics",
    tag_levels = "a"  # (a),(b)...
  ) +
  plot_layout(
    heights = c(1, 1.5)  # bottom
                          # panel wider
  )

# Alternative: cowplot
# library(cowplot)
# plot_grid(p1, p2, p3, p4,
#   ncol = 2, labels = "AUTO")
\end{lstlisting}
    \end{column}
    \begin{column}{0.45\textwidth}
      \textbf{patchwork operators:}

      \begin{itemize}
        \item \texttt{p1 | p2}: side by side
        \item \texttt{p1 / p2}: stacked
        \item \texttt{(p1 | p2) / p3}: row + row
        \item \texttt{plot\_layout(widths=)}: control
              relative panel sizes
        \item \texttt{tag\_levels}: auto-label
              panels (a), (b), (c)\ldots
      \end{itemize}
    \end{column}
  \end{columns}
\end{frame}

% ================================================================
\begin{frame}[fragile]{Multi-Panel in Python: GridSpec}
  \begin{columns}[T]
    \begin{column}{0.52\textwidth}
\begin{lstlisting}[style=pythonstyle]
from matplotlib.gridspec import (
    GridSpec)

fig = plt.figure(figsize=(8, 5))
gs = GridSpec(2, 3, figure=fig,
    hspace=0.35, wspace=0.3)

# Top row: 3 equal panels
ax1 = fig.add_subplot(gs[0, 0])
ax2 = fig.add_subplot(gs[0, 1])
ax3 = fig.add_subplot(gs[0, 2])

# Bottom: full-width panel
ax4 = fig.add_subplot(gs[1, :])

# Panel labels
for i, ax in enumerate(
    [ax1, ax2, ax3, ax4]):
    ax.text(0.02, 0.95,
        f"({chr(97+i)})",
        transform=ax.transAxes,
        fontsize=9,
        fontweight="bold",
        va="top")
\end{lstlisting}
    \end{column}
    \begin{column}{0.45\textwidth}
      \textbf{GridSpec features:}

      \begin{itemize}
        \item Arbitrary row/column spans via
              slice notation (\texttt{gs[1, :]})
        \item \texttt{hspace}, \texttt{wspace}
              control inter-panel margins
        \item \texttt{height\_ratios},
              \texttt{width\_ratios} for unequal
              panels
        \item \texttt{fig.add\_subplot()} places
              axes on the grid
      \end{itemize}
    \end{column}
  \end{columns}
\end{frame}

% ================================================================
\section{Aspect Ratio \& Export}
% ================================================================

\begin{frame}{Aspect Ratio: Banking to 45 Degrees}
  \begin{center}
    \includegraphics[width=0.92\textwidth]{figures_m05/aspect_ratio}
  \end{center}
  \begin{exampleblock}{Data Insight}
    Cleveland (1993) showed that line chart slopes are most accurately
    perceived when the average absolute slope is $\sim$45\textdegree.
    A square aspect ratio compresses temporal trends; a wider panel
    ``banks'' the slopes into the readable range.
  \end{exampleblock}
\end{frame}

% ================================================================
\begin{frame}{Export Format Guide}
  \begin{center}
    \includegraphics[width=0.85\textwidth]{figures_m05/export_formats}
  \end{center}
  \begin{alertblock}{Pro-Tip}
    Always export charts as \textbf{PDF or SVG} (vector) for
    publications and presentations. Use PNG at $\geq$ 150 dpi only when
    vector is not supported. Never use JPEG for charts --- lossy
    compression creates artefacts around sharp edges and text.
  \end{alertblock}
\end{frame}

% ================================================================
\begin{frame}[fragile]{Export in R}
  \begin{columns}[T]
    \begin{column}{0.48\textwidth}
\begin{lstlisting}[style=rstyle]
# Vector (publication)
ggsave("chart.pdf",
  plot = p,
  width = 7,   # inches
  height = 5,
  device = cairo_pdf)

# Vector (web)
ggsave("chart.svg",
  plot = p,
  width = 7, height = 5)

# Raster (presentation)
ggsave("chart.png",
  plot = p,
  width = 7, height = 5,
  dpi = 300)
\end{lstlisting}
    \end{column}
    \begin{column}{0.48\textwidth}
\begin{lstlisting}[style=pythonstyle]
# Vector (publication)
fig.savefig("chart.pdf",
    bbox_inches="tight")

# Vector (web)
fig.savefig("chart.svg",
    bbox_inches="tight")

# Raster (presentation)
fig.savefig("chart.png",
    dpi=300,
    bbox_inches="tight")

# NEVER: fig.savefig("chart.jpg")
# JPEG = lossy artefacts
\end{lstlisting}
    \end{column}
  \end{columns}
\end{frame}

% ================================================================
\section{Synthesis}
% ================================================================

\begin{frame}{Key Takeaways}
  \begin{enumerate}
    \item \textbf{Typographic hierarchy} (title $>$ subtitle $>$ body
          $>$ annotation $>$ caption) guides the reader's eye through
          the chart. Maintain a 1.2$\times$ size ratio between levels.
    \item \textbf{Sans-serif fonts} are the default for charts.
          Use \texttt{base\_family} (R) or \texttt{rcParams} (Python).
    \item \textbf{Whitespace} is structural, not decorative. Generous
          margins reduce cognitive load.
    \item \textbf{Direct labels} at data points eliminate the legend-shuttle
          problem. Use \texttt{ggrepel} / \texttt{adjustText}.
    \item \textbf{Multi-panel composition}: \texttt{patchwork} (R) and
          \texttt{GridSpec} (Python) support arbitrary layouts.
    \item \textbf{Export as vector} (PDF/SVG) for print; PNG at $\geq$ 300
          dpi for raster contexts. Never JPEG for charts.
  \end{enumerate}
\end{frame}

% ================================================================
\begin{frame}{Essential Readings for This Module}
  \begin{itemize}
    \item Butterick, M. (2015). \emph{Practical Typography}.
          \texttt{https://practicaltypography.com}
    \item Cleveland, W.~S. (1993). ``Banking to 45\textdegree.''
          In: \emph{Visualizing Data}, Hobart Press.
    \item Schwabish, J. (2021). \emph{Better Data Visualizations},
          Chapters 3--4 (Layout and Labelling).
    \item Wilke, C.~O. (2019). \emph{Fundamentals of Data Visualization},
          Chapters 22--23 (Titles, Captions, Tables).
  \end{itemize}

  \vspace{6pt}
  \textbf{Next Module}: The Design Process (W01-M06)
\end{frame}

\end{document}
